\documentclass{article}
\usepackage{graphicx}
\usepackage[T2A]{fontenc}
\usepackage[utf8]{inputenc}
\usepackage[english, russian]{babel}
\usepackage[left=30mm, right=10mm, top=20mm, bottom=20mm]{geometry}
\usepackage{framed}

\usepackage{fancyhdr}
\usepackage{amsmath}
\usepackage{pgfplots}

\pgfplotsset{compat=1.18}

\pagestyle{fancy}
\fancyhf{}
\fancyhead[R]{Шильдт Даниэль гр. 5030102/30202}
\fancyfoot[C]{\thepage}
\begin{document}
\begin{framed}
    \textbf{I.} Рассмотреть зависимость $y=1+3x_1-2x_2+x_3+e$, где случайная величина $e$ распределена
    по нормальному закону $N(0,1)$. Сгенерировать 200 различных замеров вектора факторов
    $(x_{1k},x_{2k},x_{3k})$ и вычислить соответствующие им значения отклика $y_k$. По сгенерированным
    данным построить множественную линейную регрессию.
\end{framed}
    \textbf{1.} Значения факторов $x_{1k},x_{2k},x_{3k}$ и ошибки $e_k$ сгенерированы независимо из $N(0,1)$, $k=1,\dots,n$, $n=200$. Отклик вычислен по формуле $y_k=1+3x_{1k}-2x_{2k}+x_{3k}+e_k$. Первые пять замеров:
    \begin{center}
    \begin{tabular}{|c|c|c|c|c|c|}
        \hline
        $k$ & $x_{1k}$ & $x_{2k}$ & $x_{3k}$ & $y_k$ & $\hat y_k$ \\ \hline
        1 & 0.305 & $-1.040$ & 0.750 & 5.260 & 4.688 \\
        2 & 0.941 & $-1.951$ & $-1.302$ & 5.844 & 6.501 \\
        3 & 0.128 & $-0.316$ & $-0.017$ & 3.274 & 1.970 \\
        4 & $-0.853$ & 0.879 & 0.778 & $-3.168$ & $-2.655$ \\
        5 & 0.066 & 1.127 & 0.468 & $-1.225$ & $-0.656$ \\ \hline
    \end{tabular}
    \end{center}

    \textbf{2.} Модель множественной линейной регрессии с $p=3$ факторами:
    \[y_k=\beta_0+\beta_1x_{1k}+\beta_2x_{2k}+\beta_3x_{3k}+e_k,\qquad k=1,\dots,n.\]
    В матричной форме $X=(1\ x_1\ x_2\ x_3)$ --- матрица размера $200\times4$, первый столбец которой состоит из единиц. МНК-оценки:
    \[\hat\beta=(X^TX)^{-1}X^Ty.\]
    Стандартные ошибки оценок берутся из диагонали матрицы $\hat V(\hat\beta)=s_e^2(X^TX)^{-1}$, где $s_e^2=\frac{RSS}{n-p-1}$. Для проверки значимости каждого коэффициента используется $t=\hat\beta_j/SE(\hat\beta_j)$ с $n-p-1=196$ степенями свободы.

    \textbf{3.} Результаты:
    \begin{center}
    \begin{tabular}{|c|c|c|c|c|c|}
        \hline
        Коэффициент & Истинное значение & $\hat\beta_j$ & $SE(\hat\beta_j)$ & $t$ & $p$-значение \\ \hline
        $\beta_0$ & 1 & 0.9629 & 0.0726 & 13.27 & $4\cdot10^{-29}$ \\
        $\beta_1$ & 3 & 3.0368 & 0.0722 & 42.07 & $5\cdot10^{-100}$ \\
        $\beta_2$ & $-2$ & $-2.0076$ & 0.0749 & $-26.80$ & $2\cdot10^{-67}$ \\
        $\beta_3$ & 1 & 0.9483 & 0.0776 & 12.22 & $7\cdot10^{-26}$ \\ \hline
    \end{tabular}
    \end{center}
    Уравнение регрессии:
    \[\hat y=0.963+3.037\,x_1-2.008\,x_2+0.948\,x_3.\]
    Все оценки отличаются от истинных значений меньше чем на две стандартные ошибки, то есть истинные коэффициенты попадают в 95\%-доверительные интервалы $\hat\beta_j\pm2\,SE(\hat\beta_j)$.

\begin{framed}
    \textbf{II.} Вычислить сумму квадратов остатков (RSS), стандартную ошибку остатков (RSE)
    и корреляционное отношение $\eta^2$ ($R^2$). Сделать выводы о степени адекватности линейной модели.
\end{framed}
    \textbf{1.} Сумма квадратов остатков и общая сумма квадратов:
    \[RSS=\sum_{k=1}^{n}(y_k-\hat y_k)^2\approx205.27,\qquad TSS=\sum_{k=1}^{n}(y_k-\bar y)^2\approx2732.46.\]

    \textbf{2.} Стандартная ошибка остатков:
    \[RSE=\sqrt{\frac{RSS}{n-p-1}}=\sqrt{\frac{205.27}{196}}\approx1.023.\]

    \textbf{3.} Корреляционное отношение:
    \[\eta^2=R^2=\frac{TSS-RSS}{TSS}=\frac{2732.46-205.27}{2732.46}\approx0.925.\]

    \textbf{4.} Проверим гипотезу $H_0:\beta_1=\beta_2=\beta_3=0$ с помощью $F$-статистики:
    \[F=\frac{(TSS-RSS)/p}{RSS/(n-p-1)}=\frac{(2732.46-205.27)/3}{205.27/196}\approx804.3.\]
    Критическое значение распределения Фишера $F_{0.95}(3,196)\approx2.65$, $p$-значение порядка $10^{-110}$. Гипотеза $H_0$ отвергается: отклик связан с факторами.

    \begin{center}
    \begin{tikzpicture}
    \begin{axis}[
        width=0.48\textwidth, height=6cm,
        xlabel={прогноз $\hat y$}, ylabel={наблюдаемое $y$},
        xmin=-10, xmax=12, ymin=-10, ymax=12,
        grid=both, grid style={line width=0.1pt, draw=gray!30},
        title={Наблюдаемые и предсказанные значения}]
        \addplot[only marks, mark=*, mark size=1pt, blue] coordinates {
            (4.69,5.26) (6.50,5.84) (1.97,3.27) (-2.66,-3.17) (-0.66,-1.23) (-3.30,-2.73) (3.56,4.31) (-3.71,-3.19) (0.87,-0.73) (1.65,2.24) (7.80,6.69) (-1.67,-1.31) (1.52,0.09) (1.96,2.45) (-1.41,-2.15) (0.09,-1.15) (3.16,3.90) (5.50,5.73) (0.56,-0.02) (8.16,9.64) (-3.32,-3.66) (2.07,2.16) (1.65,2.01) (1.87,1.22) (-2.69,-2.14) (4.07,3.55) (3.92,3.54) (-0.68,0.72) (-4.09,-5.07) (0.46,-1.89) (2.78,4.38) (2.74,5.24) (-0.67,-0.98) (0.20,-1.60) (-2.48,-2.79) (1.92,1.66) (7.39,7.22) (1.51,0.34) (2.47,3.09) (-2.41,-1.85) (-2.30,-3.70) (3.75,4.46) (7.19,9.18) (-2.37,-2.20) (1.94,1.54) (-1.96,-1.98) (1.57,2.19) (2.49,0.83) (-1.25,-1.06) (-1.46,-1.79) (0.97,2.77) (1.27,1.24) (-1.09,0.59) (-1.39,-2.45) (9.87,10.45) (-3.28,-2.87) (-0.05,-0.91) (2.07,2.20) (-1.77,-0.39) (1.96,1.68) (4.10,2.38) (4.92,4.98) (2.68,1.73) (2.35,1.98) (-4.18,-4.26) (-3.42,-5.07) (1.60,0.03) (5.67,6.17) (-5.96,-6.42) (-1.29,-3.80) (6.44,7.22) (1.10,1.48) (-0.23,-0.93) (-0.47,-1.71) (3.76,4.62) (-1.62,-1.12) (-2.44,-0.07) (-0.32,0.22) (-1.97,-1.64) (2.10,1.94) (-3.18,-2.38) (-8.99,-8.32) (-3.30,-1.81) (-1.53,-1.96) (7.96,7.48) (2.09,1.61) (3.09,3.90) (-0.84,0.68) (0.51,0.09) (1.92,1.52) (-1.96,-1.58) (4.53,4.27) (2.01,3.08) (-0.86,-3.17) (1.81,0.94) (-0.95,-1.67) (0.04,-2.24) (-3.96,-5.02) (0.17,-0.68) (3.67,3.52) (10.11,11.23) (-0.91,-1.16) (-2.11,-3.05) (4.28,4.22) (6.16,7.23) (-0.73,-1.70) (7.08,6.20) (-1.26,-0.61) (3.53,3.21) (2.62,3.27) (3.61,3.55) (-3.34,-2.62) (-0.61,-1.52) (7.70,7.70) (1.88,1.79) (-3.56,-4.63) (-6.45,-7.89) (5.45,6.18) (3.96,3.62) (3.38,2.37) (0.68,0.68) (-0.60,0.68) (0.55,-1.66) (4.76,3.34) (2.40,1.52) (6.54,7.85) (6.86,7.12) (6.97,7.65) (-1.89,-2.92) (0.66,-0.51) (-6.61,-6.06) (1.08,1.17) (3.31,3.90) (-0.63,-0.75) (4.74,5.05) (1.18,1.43) (-3.89,-3.09) (-4.05,-3.09) (0.73,-0.78) (6.42,4.23) (2.36,3.36) (-2.09,-0.84) (-1.36,-2.43) (4.70,2.87) (5.12,5.30) (-0.82,0.11) (-2.06,-0.21) (2.49,3.10) (3.90,3.62) (-1.92,-2.73) (1.71,1.87) (7.28,6.95) (6.24,5.24) (-0.54,1.51) (-0.76,1.05) (2.00,3.28) (6.58,5.67) (0.47,1.39) (4.53,5.24) (-0.05,0.60) (-1.15,0.14) (2.68,3.32) (-3.08,-2.31) (1.15,1.83) (-1.11,-0.73) (7.16,5.35) (-0.79,-0.53) (3.70,3.07) (-4.65,-5.67) (3.85,5.57) (-3.09,-1.26) (-0.56,0.86) (5.23,5.53) (7.02,8.29) (6.73,6.81) (-3.18,-2.90) (-1.74,-2.81) (-0.57,-0.08) (-5.58,-5.45) (8.51,7.14) (-2.26,-2.21) (0.78,0.73) (4.28,6.13) (2.73,3.67) (-2.31,-2.26) (0.56,0.80) (-1.10,-0.98) (1.15,0.93) (2.06,0.98) (-1.94,-2.02) (-3.05,-3.69) (-3.53,-4.78) (2.28,2.77) (5.81,6.15) (2.90,1.19) (8.60,8.38) (-2.72,-1.74) (-3.50,-2.49) (-2.57,-1.49) (1.26,1.22)};
        \addplot[red, thick, domain=-10:12] {x};
    \end{axis}
    \end{tikzpicture}
    \hfill
    \begin{tikzpicture}
    \begin{axis}[
        width=0.48\textwidth, height=6cm,
        xlabel={прогноз $\hat y$}, ylabel={остаток $y-\hat y$},
        xmin=-10, xmax=12, ymin=-3.5, ymax=3.5,
        grid=both, grid style={line width=0.1pt, draw=gray!30},
        title={Остатки}]
        \addplot[only marks, mark=*, mark size=1pt, blue] coordinates {
            (4.69,0.57) (6.50,-0.66) (1.97,1.30) (-2.66,-0.51) (-0.66,-0.57) (-3.30,0.56) (3.56,0.76) (-3.71,0.51) (0.87,-1.61) (1.65,0.59) (7.80,-1.11) (-1.67,0.37) (1.52,-1.43) (1.96,0.49) (-1.41,-0.74) (0.09,-1.24) (3.16,0.74) (5.50,0.23) (0.56,-0.58) (8.16,1.48) (-3.32,-0.34) (2.07,0.09) (1.65,0.36) (1.87,-0.64) (-2.69,0.55) (4.07,-0.52) (3.92,-0.38) (-0.68,1.40) (-4.09,-0.98) (0.46,-2.35) (2.78,1.60) (2.74,2.50) (-0.67,-0.31) (0.20,-1.80) (-2.48,-0.30) (1.92,-0.27) (7.39,-0.17) (1.51,-1.16) (2.47,0.62) (-2.41,0.56) (-2.30,-1.41) (3.75,0.71) (7.19,1.99) (-2.37,0.17) (1.94,-0.40) (-1.96,-0.02) (1.57,0.62) (2.49,-1.66) (-1.25,0.20) (-1.46,-0.33) (0.97,1.80) (1.27,-0.03) (-1.09,1.68) (-1.39,-1.07) (9.87,0.58) (-3.28,0.42) (-0.05,-0.86) (2.07,0.13) (-1.77,1.38) (1.96,-0.28) (4.10,-1.72) (4.92,0.06) (2.68,-0.95) (2.35,-0.36) (-4.18,-0.08) (-3.42,-1.65) (1.60,-1.56) (5.67,0.50) (-5.96,-0.45) (-1.29,-2.51) (6.44,0.78) (1.10,0.38) (-0.23,-0.69) (-0.47,-1.24) (3.76,0.86) (-1.62,0.49) (-2.44,2.37) (-0.32,0.54) (-1.97,0.34) (2.10,-0.16) (-3.18,0.79) (-8.99,0.67) (-3.30,1.49) (-1.53,-0.43) (7.96,-0.48) (2.09,-0.49) (3.09,0.81) (-0.84,1.52) (0.51,-0.42) (1.92,-0.40) (-1.96,0.38) (4.53,-0.25) (2.01,1.07) (-0.86,-2.31) (1.81,-0.88) (-0.95,-0.72) (0.04,-2.29) (-3.96,-1.05) (0.17,-0.85) (3.67,-0.15) (10.11,1.12) (-0.91,-0.25) (-2.11,-0.94) (4.28,-0.05) (6.16,1.07) (-0.73,-0.96) (7.08,-0.88) (-1.26,0.65) (3.53,-0.33) (2.62,0.65) (3.61,-0.06) (-3.34,0.71) (-0.61,-0.91) (7.70,0.00) (1.88,-0.09) (-3.56,-1.08) (-6.45,-1.44) (5.45,0.73) (3.96,-0.34) (3.38,-1.01) (0.68,0.00) (-0.60,1.28) (0.55,-2.21) (4.76,-1.42) (2.40,-0.88) (6.54,1.31) (6.86,0.26) (6.97,0.68) (-1.89,-1.03) (0.66,-1.17) (-6.61,0.55) (1.08,0.09) (3.31,0.59) (-0.63,-0.12) (4.74,0.31) (1.18,0.25) (-3.89,0.80) (-4.05,0.96) (0.73,-1.51) (6.42,-2.19) (2.36,1.00) (-2.09,1.25) (-1.36,-1.07) (4.70,-1.82) (5.12,0.17) (-0.82,0.93) (-2.06,1.85) (2.49,0.61) (3.90,-0.28) (-1.92,-0.82) (1.71,0.16) (7.28,-0.33) (6.24,-1.01) (-0.54,2.05) (-0.76,1.81) (2.00,1.29) (6.58,-0.91) (0.47,0.91) (4.53,0.71) (-0.05,0.66) (-1.15,1.29) (2.68,0.64) (-3.08,0.77) (1.15,0.68) (-1.11,0.38) (7.16,-1.81) (-0.79,0.25) (3.70,-0.64) (-4.65,-1.01) (3.85,1.72) (-3.09,1.83) (-0.56,1.42) (5.23,0.31) (7.02,1.27) (6.73,0.07) (-3.18,0.27) (-1.74,-1.07) (-0.57,0.49) (-5.58,0.13) (8.51,-1.37) (-2.26,0.05) (0.78,-0.05) (4.28,1.86) (2.73,0.94) (-2.31,0.05) (0.56,0.24) (-1.10,0.12) (1.15,-0.22) (2.06,-1.08) (-1.94,-0.08) (-3.05,-0.64) (-3.53,-1.25) (2.28,0.50) (5.81,0.34) (2.90,-1.71) (8.60,-0.21) (-2.72,0.97) (-3.50,1.02) (-2.57,1.08) (1.26,-0.04)};
        \addplot[black, domain=-10:12] {0};
    \end{axis}
    \end{tikzpicture}
    \end{center}

    \textbf{Вывод:} модель объясняет около 92.5\% разброса отклика, $F$-критерий подтверждает наличие связи, все коэффициенты значимы и близки к истинным. $RSE\approx1.02$ практически совпадает со стандартным отклонением ошибки $\sigma=1$, то есть оставшийся разброс объясняется только случайной ошибкой $e$. Остатки расположены вокруг нуля без видимой закономерности. Линейная модель адекватна.

\begin{framed}
    \textbf{III.} На сайте http://www.pogodaiklimat.ru/file.htm найти данные о среднегодовых температурах
    за период наблюдений. Построить полиномиальные линейные регрессии температуры от времени.
    Реализовать процедуру прямого отбора при выборе степени полинома.
\end{framed}
    \textbf{1.} Взяты среднегодовые температуры Тюмени за последние 40 лет наблюдений (1986--2025 гг.):
    \begin{center}
    \small
    \begin{tabular}{|c|c|c|c|c|c|c|c|c|c|c|}
        \hline
        Год & 1986 & 1987 & 1988 & 1989 & 1990 & 1991 & 1992 & 1993 & 1994 & 1995 \\ \hline
        $T$, °C & 0.6 & 1.2 & 2.7 & 2.5 & 2.9 & 2.9 & 1.6 & 1.2 & 1.5 & 4.1 \\ \hline
        \hline
        Год & 1996 & 1997 & 1998 & 1999 & 2000 & 2001 & 2002 & 2003 & 2004 & 2005 \\ \hline
        $T$, °C & 1.3 & 2.0 & 1.4 & 2.0 & 2.4 & 1.6 & 2.1 & 3.3 & 2.7 & 3.1 \\ \hline
        \hline
        Год & 2006 & 2007 & 2008 & 2009 & 2010 & 2011 & 2012 & 2013 & 2014 & 2015 \\ \hline
        $T$, °C & 1.8 & 3.3 & 3.7 & 2.0 & 1.4 & 2.0 & 2.7 & 3.0 & 1.7 & 3.1 \\ \hline
        \hline
        Год & 2016 & 2017 & 2018 & 2019 & 2020 & 2021 & 2022 & 2023 & 2024 & 2025 \\ \hline
        $T$, °C & 2.8 & 2.6 & 1.6 & 2.8 & 4.7 & 2.9 & 2.9 & 4.3 & 3.6 & 4.4 \\ \hline
    \end{tabular}
    \end{center}

    \textbf{2.} Полиномиальная регрессия степени $p$:
    \[T=\beta_0+\beta_1z+\beta_2z^2+\dots+\beta_pz^p+e,\]
    сводится к множественной линейной регрессии введением факторов $X_j=z^j$, $j=1,\dots,p$. Здесь $z_k=(t_k-\bar t)/s_t$ --- год, приведённый к нулевому среднему и единичному разбросу ($\bar t=2005.5$, $s_t\approx11.54$). Без этого значения $t_k^p$ при $t_k\approx2000$ отличаются на много порядков и матрица $X^TX$ получается плохо обусловленной.

    \textbf{3.} Прямой отбор: начинаем с нулевой модели $T=\beta_0$ и на каждом шаге добавляем следующую степень $z^p$. Новая степень остаётся в модели, если её коэффициент значим ($p$-значение $t$-статистики меньше 0.05). Как только добавленная степень оказывается незначимой, процедура останавливается. Характеристики моделей степеней от 1 до 6:
    \begin{center}
    \begin{tabular}{|c|c|c|c|c|c|}
        \hline
        Степень $p$ & $RSS$ & $RSE$ & $R^2$ & $t$ нового коэффициента & $p$-значение \\ \hline
        1 & 26.897 & 0.841 & 0.267 & 3.72 & 0.001 \\
        2 & 25.967 & 0.838 & 0.292 & 1.15 & 0.257 \\
        3 & 24.042 & 0.817 & 0.345 & 1.70 & 0.098 \\
        4 & 24.017 & 0.828 & 0.345 & 0.19 & 0.851 \\
        5 & 22.922 & 0.821 & 0.375 & 1.27 & 0.211 \\
        6 & 20.181 & 0.782 & 0.450 & $-2.12$ & 0.042 \\ \hline
    \end{tabular}
    \end{center}
    На первом шаге коэффициент при $z$ значим ($p=0.001$), на втором коэффициент при $z^2$ незначим ($p=0.257$). Процедура останавливается, выбрана степень $p=1$.

    \textbf{4.} Выбранная модель:
    \[\hat T=2.510+0.494\,z=2.510+0.0428\,(t-2005.5),\]
    то есть среднегодовая температура растёт примерно на 0.43~°C за 10 лет. Для неё $RSE\approx0.84$~°C, $R^2\approx0.27$, $F\approx13.8$ ($p$-значение 0.0006).

    \begin{center}
    \begin{tikzpicture}
    \begin{axis}[
        width=0.9\textwidth, height=7cm,
        xlabel={год}, ylabel={температура, °C},
        xmin=1984, xmax=2027, ymin=0, ymax=5,
        x tick label style={/pgf/number format/1000 sep={}},
        grid=both, grid style={line width=0.1pt, draw=gray!30},
        legend pos=north west, legend cell align=left]
        \addplot[only marks, mark=*, mark size=1.5pt, gray] coordinates {
            (1986,0.6) (1987,1.2) (1988,2.7) (1989,2.5) (1990,2.9) (1991,2.9) (1992,1.6) (1993,1.2) (1994,1.5) (1995,4.1) (1996,1.3) (1997,2.0) (1998,1.4) (1999,2.0) (2000,2.4) (2001,1.6) (2002,2.1) (2003,3.3) (2004,2.7) (2005,3.1) (2006,1.8) (2007,3.3) (2008,3.7) (2009,2.0) (2010,1.4) (2011,2.0) (2012,2.7) (2013,3.0) (2014,1.7) (2015,3.1) (2016,2.8) (2017,2.6) (2018,1.6) (2019,2.8) (2020,4.7) (2021,2.9) (2022,2.9) (2023,4.3) (2024,3.6) (2025,4.4)};
        \addlegendentry{среднегодовая температура}
        \addplot[blue, very thick] coordinates {
            (1986,1.675) (1987,1.718) (1988,1.760) (1989,1.803) (1990,1.846) (1991,1.889) (1992,1.932) (1993,1.975) (1994,2.017) (1995,2.060) (1996,2.103) (1997,2.146) (1998,2.189) (1999,2.232) (2000,2.274) (2001,2.317) (2002,2.360) (2003,2.403) (2004,2.446) (2005,2.489) (2006,2.531) (2007,2.574) (2008,2.617) (2009,2.660) (2010,2.703) (2011,2.746) (2012,2.788) (2013,2.831) (2014,2.874) (2015,2.917) (2016,2.960) (2017,3.003) (2018,3.045) (2019,3.088) (2020,3.131) (2021,3.174) (2022,3.217) (2023,3.260) (2024,3.302) (2025,3.345)};
        \addlegendentry{степень 1 (выбрана)}
        \addplot[green!60!black, dashed] coordinates {
            (1986,1.991) (1987,1.985) (1988,1.982) (1989,1.981) (1990,1.983) (1991,1.988) (1992,1.995) (1993,2.004) (1994,2.016) (1995,2.031) (1996,2.048) (1997,2.068) (1998,2.090) (1999,2.115) (2000,2.142) (2001,2.173) (2002,2.205) (2003,2.240) (2004,2.278) (2005,2.318) (2006,2.361) (2007,2.406) (2008,2.454) (2009,2.505) (2010,2.558) (2011,2.614) (2012,2.672) (2013,2.733) (2014,2.796) (2015,2.862) (2016,2.930) (2017,3.001) (2018,3.075) (2019,3.151) (2020,3.230) (2021,3.311) (2022,3.395) (2023,3.481) (2024,3.570) (2025,3.662)};
        \addlegendentry{степень 2}
        \addplot[orange, dashed] coordinates {
            (1986,1.492) (1987,1.640) (1988,1.770) (1989,1.883) (1990,1.981) (1991,2.065) (1992,2.136) (1993,2.194) (1994,2.241) (1995,2.278) (1996,2.307) (1997,2.327) (1998,2.341) (1999,2.349) (2000,2.352) (2001,2.352) (2002,2.350) (2003,2.347) (2004,2.343) (2005,2.340) (2006,2.339) (2007,2.342) (2008,2.348) (2009,2.360) (2010,2.378) (2011,2.404) (2012,2.438) (2013,2.482) (2014,2.537) (2015,2.603) (2016,2.683) (2017,2.776) (2018,2.885) (2019,3.010) (2020,3.152) (2021,3.313) (2022,3.493) (2023,3.693) (2024,3.916) (2025,4.161)};
        \addlegendentry{степень 3}
    \end{axis}
    \end{tikzpicture}
    \end{center}

    \textbf{Вывод:} $R^2$ растёт с каждой новой степенью (от 0.267 до 0.450), потому что полином большей степени всегда лучше подгоняется к тем же точкам. Однако уже вторая степень не даёт значимого улучшения, и прямой отбор выбирает линейную модель. У степени 6 $p$-значение формально меньше 0.05, но до неё процедура не доходит; при 40 точках полином шестой степени в основном подстраивается под случайные колебания отдельных лет. Линейная модель показывает значимый рост температуры, но объясняет только около 27\% разброса: остальное приходится на колебания от года к году, которые трендом не описываются.

\end{document}
